% Generated draft: requires mathematical and source review.
\documentclass[11pt]{article}
\usepackage[margin=25mm]{geometry}
\usepackage{fontspec}
\setmainfont{Latin Modern Roman}
\usepackage{amsmath,amssymb,mathtools,unicode-math}
\setmathfont{Latin Modern Math}
\usepackage{enumitem,needspace,xurl,hyperref}
\setlength{\parindent}{0pt}
\setlength{\parskip}{5pt}
\setlength{\emergencystretch}{3em}
\newcommand{\N}{\mathbb{N}}
\newcommand{\Z}{\mathbb{Z}}
\newcommand{\Q}{\mathbb{Q}}
\newcommand{\R}{\mathbb{R}}
\newcommand{\C}{\mathbb{C}}
\newcommand{\F}{\mathbb{F}}
\newcommand{\A}{\mathbb{A}}
\newcommand{\PP}{\mathbb{P}}
\newcommand{\E}{\mathbb{E}}
\newcommand{\eps}{\varepsilon}
\newcommand{\dd}{\,\mathrm{d}}
\begin{document}
\section*{UnsolvedMath ID: 700544}
\textbf{Source problem code:} catalogue entry
\subsection{Definitions and mathematical statement}
Let \(G\) be a group and let \(S=S^{-1}\subseteq G\) be a finite generating set. Define the word length \(\ell_S(g)\) to be the least number of elements of \(S\) whose product is \(g\), with \(\ell_S(1_G)=0\). Define the ball growth function by
\[
\gamma_{G,S}(n)=\bigl|\{g\in G:\ell_S(g)\leq n\}\bigr|,
\qquad n\in\Z_{\geq 0}.
\]
For subgroups \(A,B\leq G\), let \([A,B]\) denote the subgroup generated by the commutators \(a^{-1}b^{-1}ab\). A group \(N\) is nilpotent if its lower central series, defined by \(N_1=N\) and \(N_{j+1}=[N_j,N]\), satisfies \(N_{r+1}=\{1_G\}\) for some integer \(r\geq 1\). A group \(G\) is nilpotent-by-finite if it has a nilpotent normal subgroup \(N\trianglelefteq G\) with \([G:N]<\infty\).

The question attributed to R. I. Grigorchuk is whether there exist a group \(G\), a finite symmetric generating set \(S\) of \(G\), and a real constant \(c>1\) such that \(G\) is not nilpotent-by-finite and
\[
\forall n\in\Z_{\geq 0},\qquad \gamma_{G,S}(n)\leq c^{\sqrt n}.
\]
The source does not explicitly define its growth function or specify whether the majorization is required for every argument or only eventually. The formulation above uses word-ball growth and an everywhere upper bound. For this growth convention, an eventual upper bound of the stated form can be made an everywhere upper bound by increasing \(c\). The existence of such a bound is also independent of the finite generating set, after changing \(c\). The intended growth convention still requires source confirmation. The registry labels this problem open; this is supplied source metadata, not an independently verified status assessment.
\subsection{Short English statement}
Does there exist a finitely generated group with no nilpotent normal subgroup of finite index whose word-ball growth is bounded above by \(c^{\sqrt n}\) for some constant \(c>1\)? This uses the stated growth convention and requires human review against the source.
\subsection{Sources}
\begin{itemize}
\item Local UnsolvedMath registry, permanent ID \(700544\), exact source problem code \texttt{catalogue entry}, title \emph{Kourovka Notebook Problem 9.9}. Supplied upstream revision: \texttt{d371ee917ac9647fbb06cd6fc529deec6f3dda82}. The record was created and updated on \texttt{2026-10-16T14:00:00Z} and carries the supplied status \texttt{open}.
\item \emph{The Kourovka Notebook (Unsolved Problems in Group Theory)}, 9th issue (1984), Problem 9.9; proposed by R. I. Grigorchuk. The supplied background also identifies the 21st issue edition (2026), arXiv:1401.0300v48. These bibliographic details are supplied attribution and have not been independently verified.
\item The supplied collection description attributes editing of the Kourovka Notebook, Issues 1--21 (1965--2026), to E. I. Khukhro and V. D. Mazurov.
\end{itemize}
\end{document}
